\documentclass[11pt]{article}
\usepackage[]{authblk}
\usepackage{graphicx}
\usepackage{hanging}
\usepackage{indentfirst}
\usepackage{fancyhdr}
\usepackage{setspace}
\usepackage{multicol}
\usepackage{color}
\usepackage{lipsum}
\usepackage[breaklinks]{hyperref}
\usepackage[all]{hypcap}
\definecolor{darkblue}{rgb}{0.0,0.0,0.5}
%\onehalfspace
\setlength{\evensidemargin}{0in}
\setlength{\oddsidemargin}{0in}
\setlength{\paperheight}{11in}
\setlength{\paperwidth}{8.5in}
\setlength{\tabcolsep}{0in}
\setlength{\textheight}{8.5in} % This changed from 9in
\setlength{\textwidth}{6.5in}
\setlength{\topmargin}{0in}
\setlength{\topskip}{0mm}
\setlength{\voffset}{-0.5in}
\setlength{\parindent}{1cm}
%\setlength{\headheight}{99pt}
\setlength{\headheight}{0.5in}
\parskip = 0mm
\pagestyle{plain}

\definecolor{citescol}{RGB}{73,0,165}
\definecolor{urlscol}{RGB}{3,0,255}
\definecolor{linkscol}{RGB}{187,24,0}
\definecolor{isured}{RGB}{206,1,38}
\definecolor{isukhaki}{RGB}{173,160,122}

\hypersetup{colorlinks=true,linkcolor=linkscol,citecolor=citescol,urlcolor=urlscol}


\renewcommand{\headrulewidth}{0pt}
%\renewcommand{\footrulewidth}{0.0pt}
\lhead{{\textcolor{isured}{\LARGE{\textsc{Iowa State University}}}}\\{\textcolor{isukhaki}{OF SCIENCE AND TECHNOLOGY}}}
\rhead{\scriptsize{\includegraphics[width=0.85in]{assets/EEOB.eps}\\{\textcolor{isured}{Ecology, Evolution, and Organismal Biology}}}}

\begin{document}

\thispagestyle{fancy}
\begin{flushright}
\today
\end{flushright}
\vspace{2mm}

%\noindent Chair, Macroevolution Search Committee\\
%Department of Ecology, Evolution, and Organismal Biology\\
%Iowa State University\\
%Ames, Iowa 50011\\


\noindent To whom it may concern,

As the lead sponsoring scientist supporting \textbf{Dr.\ April Wright}, I am writing this statement to communicate my role in her postdoctoral training and the suitability of Iowa State University (ISU) and the Field Museum for her proposed work. 
April's previous research and current interests overlap with my research program which focuses on the development of statistical methods for phylogenetic inference and macroevolution. %---with a particular emphasis on integrating fossils into modern phylogenetic analyses.

I will start by saying that April is currently a half-time postdoc in my lab (since August of 2015). 
I realize that it is unconventional for an applicant to have already joined the lab that they are proposing to work with, but April's current appointment is only part time and as the parent of a newborn, she made arrangements to start her postdoctoral research in a way that would be ideal for her and her family. 
After finishing her PhD at the University of Texas in May of 2015, April took 1.5 months off for maternity leave.
In July, she started a project with Dr.\ Mark Holder at the University of Kansas developing morphological models for inferring phylogenies (NSF DEB-1256993). 
The funds for April's work with Dr.\ Holder are only budgeted to support a postdoctoral researcher with 50\% salary and no benefits through June of 2016, thus April needed another half-time appointment. 
Notably, the nature of April's research as a computational evolutionary biologist enables her to work in two places at once. 
My lab and ISU are a good fit for April for several reasons: my research overlaps with her interests and experience, the proximity of Ames to April's family is ideal for her as a new parent, and Ames is an incredibly affordable place to live. 
For these reasons and because I know April to be a driven and excellent researcher, I recruited her to my lab as a postdoc in August of this year.
I currently have start-up funds to pay her 50\% salary and full benefits for one year.  
This arrangement has been an excellent way for April to establish herself in the field of statistical phylogenetics and within the research community at ISU. 
In my lab, I have given April the opportunity to conduct independent work and we are currently collaborating on the development of research projects that will be complementary to the work she is doing with Dr.\ Holder and the work she is proposing for the NSF fellowship.

As of October 2015, my lab is supported by start-up funds that last through the end of 2018. 
Additionally, I have two full proposals pending with NSF DEB programs. 
The first proposal, entitled ``Collaborative Research: Advancing Bayesian Phylogenetic Methods for Synthesizing Paleontological and Neontological Data'' (submitted to Phylogenetic Systematics), focuses on the development of statistical methods for incorporating stratigraphic information and biogeography into methods integrating fossil and extant taxa in Bayesian phylogenetic analyses. 
This work is in collaboration with Dr.~Robert Meredith (Montclair State University) and Dr.~Daniel Ksepka (The Bruce Museum), and will involve applying new methods to two vertebrate groups--penguins and crocodylomorphs--and relies heavily on museum collections. 
The second pending proposal, entitled ``Multilocus Analyses of Co-diversification and Phylogenetic Incongruence Between Highly Coevolved Figs and Fig Wasps'' (submitted to Evolutionary Processes), is a collaboration with Dr.\ John Nason (ISU) and Dr.\ Allen Harre (Smithsonian Tropical Research Insititute) and seeks to understand co-evolution in the fig/fig-wasp system.

A large component of my research program involves developing fully integrative Bayesian phylogenetic methods in the program RevBayes (\href{http://revbayes.github.io/about.html}{http://revbayes.github.io}). 
RevBayes is an extensive and collaboratively developed software platform for statistical phylogenetics.
April will have the opportunity to become part of the RevBayes development team and implement new methods in this open-source project. 
The work that April is proposing seeks to identify the limitations of birth-death models for estimating tree topologies and divergence times when fossil recovery rates vary across the tree. 
This builds on work I have done introducing the fossilized birth-death model for calibrating divergence-time estimates and phylogenetic inference (Heath, Huelsenbeck, \& Stadler, \textit{PNAS}, 2014; Zhang et al.\ \textit{Syst.\ Biol.}, 2015; Gavryushkina et al.\ arXiv:1506.04797, 2015). 
One of my primary collaborators on these projects is Tanja Stadler (ETH Z\"{u}rich), widely acknowledged as one of the leading experts on birth-death processes. 
Dr.~Stadler and I continue to collaborate on related projects and April will have opportunities to work with her to extend these mathematical models. 
Moreover, as a member of my lab, April will be exposed to the other areas of my research program, including the development of relaxed molecular clock models, modeling gene-family evolution, and methods for understanding coevolution.
The work April proposes and her collaborations will establish the foundation of her future work as an independent researcher.

The Department of Ecology, Evolution and Organismal Biology (EEOB) at ISU is an excellent environment for April to start her independent research career. 
The EEOB department provides a supportive and collegial community for early-career scientists.
April will interact with several leading researchers in biodiversity and systematics as an EEOB psotdoc (including Dean Adams, Anne Bronikowski, Lynn Clark, Dennis Lavrov, John Nason, Jeanne Serb, Amy Toth, Nicole Valenzuela, Jonathan Wendel). 
Research in EEOB at Iowa State is inherently organismal and offers numerous opportunities for April to apply her computational expertise to collaborations with empiricists working on important biological problems.
EEOB and ISU provide a number of training opportunities for postdoctoral researchers including career development and teaching workshops as well as opportunities to participate on academic committees and attend faculty meetings (as a non-voting observer). 
Since her arrival at ISU, April has already observed several faculty meetings and attended pedagogical workshops through the Center for Excellence in Learning and Teaching (CELT).
April has stated that she aspires to an academic career and intends to pursue a tenure-track faculty position at an institution that values undergraduate education.
Therefore, April will apply to participate in the Preparing Future Faculty (PFF) Program at ISU. PFF is a national movement sponsored by the Association of American Colleges and Universities that offers mentoring and career development instruction to early-career scientists above what is available at the departmental level. This training is designed to develop skills to enhance success in a competitive academic job market.
%In addition, t
%At Iowa State, April also has access to a thriving and interdisciplinary statistics department as well as several other life science departments (Animal Science, Agronomy, Botany, Entomology, Genetics, Development \& Cell Biology, etc.).

The work April has described in her proposal will require access to high-performance computing (HPC) resources.
ISU has an excellent facility for computationally intensive research and as a member of my lab, April has access to a newly built computing cluster in this HPC center (144-nodes $\times$ 16 cores) as well as several other computing resources.
ISU also has a wide range of programs for undergraduate research in the sciences, particularly focused on broadening the participation of underrepresented minorities, women, and first-generation college students. 
In particular, as a member of my lab, April can mentor students through the Academic Program for Excellence (\href{http://www.multicultural.dso.iastate.edu/prospective-students/apex}{http://www.multicultural.dso.iastate.edu/prospective-students/apex}) and Science Bound (\href{http://www.education.iastate.edu/science-bound.html}{http://www.education.iastate.edu/science-bound.html}).
Throughout her career, April has established a record of thoughtful and independent outreach and I'm certain she will continue this work.
These programs will contribute to her postdoctoral training by linking her with undergraduate researchers from diverse backgrounds. 
Importantly, in the short time that April has been at Iowa State, she has proven to be a valuable member of my group and the broader scientific community.
She is dedicated to helping train the next generation of computational evolutionary biologists by teaching computational skills to new EEOB graduate students through a programming study group.
April is also working with Megan O'Donnell, a member of the ISU Libraries staff, to organize a Data Carpentry course to bring open-data practices to ISU scientists.
These activities, on top of her strong record of research, demonstrate that April has the skills needed to become an active leader at an academic institution. 

The methods April will develop will be generally applicable to a diverse array of empirical studies using specimens and data from natural history collections and museum databases.
Moreover, she will have the expertise to execute such studies and apply her work to biological collections data.
April has sought out an ideal co-sponsor, Dr.~Corrie Moreau, for this component of her proposed research. 
Dr.~Moreau is the associate curator of ants at the Field Museum of Natural History in Chicago. 
Her research---supported by several grants from NSF ($>$\$1M from DEB and IOS)---focuses on ant evolution, ecology, and systematics in a broad sense with specific projects on ant-symbiont coevolution, intelligence and brain evolution in ants, and understanding the evolution of interactions between ants and their microbial environments. 
April seeks to apply her new methods to understanding macroevolution in ants.
Thus, Dr.~Moreau's expertise and access to collections at the Field Museum are essential.
Chicago and the Field Museum are easily accessible from Ames, IA and April will be able to visit frequently to take advantage of career-development opportunities and interact with a diverse group of museum-based researchers.
Since the Field Museum engages a diverse, public audience of all ages, April will have the unique opportunity to learn to communicate science outside of the academic setting.

Dr.~Moreau and I are committed to work with April to plan her training and identify activities that will prepare her for a future career as a principal investigator.
We have already discussed the research skills that she hopes to develop.
My primary objective as a supporting scientist for this work is to ensure that April builds an independent research program and gains the experience necessary for her to be competitive in a difficult academic job market.
In my lab and as a developer of RevBayes, April will strengthen her understanding of Bayesian phylogenetics, Markov chain Monte Carlo methods, mathematical modeling, and object-oriented programming.
April is a strong Python programmer and with such expertise will continue to contribute to training new personnel in my lab, activities that will hone her skills as a mentor and educator. 
RevBayes is a collaborative software project for which we have regular virtual meetings, hack-a-thons, and workshops. 
As a RevBayes developer, April will participate in all of these activities, including assisting in the teaching of workshops.
At the Field Museum, April will work with Dr.~Moreau on the empirical analyses she has outlined in her proposal.
Dr.~Moreau and her lab will educate April in ant morphology and evolution, as well as in collections-based research in general. 

April has proven that she is a collaborative and rigorous computational biologist, and is already an integral member of my lab. 
I am confident that she will become a key player in the development of statistical phylogenetic methods and in the broader evolutionary biology community through her research and outreach.

\vspace{6mm}
\noindent Sincerely,\\
%\includegraphics[width=2in]{/Users/tracyh/EEOB_Letter_Template/assets/signature2.eps}\\
\noindent Tracy A. Heath\\
Assistant Professor\\
Dept. Ecology, Evolution, \& Organismal Biology\\
Iowa State University\\
243 Bessey Hall\\
Ames, IA 50011 USA\\
\href{mailto:phylo@iastate.edu}{{phylo@iastate.edu}}\\
(515) 294-3931

\end{document}
